% GENERATED FILE. Edit canonical lessons, then run npm run generate:books.
% canonical-source-hash: fnv1a64:fb9a32fdebdf7f43
% canonical-lessons: TE-C244-vicu, TE-C244-prakasincu, TE-C244-gaddakattu, TE-C244-vadipovu, TE-C244-morugu

\chapter{To blow, To shine, To freeze, To wilt, To bark}
\label{ch:te-to-blow-to-shine-to-freeze-to-wilt-to-bark}
\hlchaptermodality{244}

\begin{chapteropening}
\textbf{By the end of this chapter:} \emph{I can say to blow, to shine, to freeze, to wilt and to bark.}

Complete the last lesson of chapter 244: I can say to blow, to shine, to freeze, to wilt and to bark.
\end{chapteropening}

\section[\texorpdfstring{\te{వీచు} (vīcu)}{వీచు (vīcu)}]{\te{వీచు} (vīcu) — to blow (of the wind)}

\label{lesson:TE-C244-vicu}



\begin{quote}
Before the new one: say the Telugu for to inform, then the Telugu for to transfer.
\end{quote}

\subsection*{You'll want to know: \te{వీచు}}
\textbf{\te{వీచు}} — \emph{vīcu} — \textquotedblleft{}to blow (of the wind)\textquotedblright{}.

\begin{grammarlens}[title={What you've built}]
One more word for this chapter.
\end{grammarlens}

\subsection*{Your turn}
\begin{itemize}
\raggedright
  \item \emph{Say it:} \emph{vīcu}
  \item \emph{Say it:} \emph{vīcu}, once more
  \item \emph{Say it:} say \emph{vīcu}
  \item \emph{Recall:} say the Telugu for to inform, then the Telugu for to transfer, then say \emph{vīcu} again
\end{itemize}

\begin{tcolorbox}[breakable,colback=teal!4,colframe=teal!35!black,title={Before you move on}]
What does \te{వీచు} mean? (\textquotedblleft{}To blow (of the wind)\textquotedblright{}.) Say it once more.
\end{tcolorbox}
\section[\texorpdfstring{\te{ప్రకాశించు} (prakāśiñcu)}{ప్రకాశించు (prakāśiñcu)}]{\te{ప్రకాశించు} (prakāśiñcu) — to shine}

\label{lesson:TE-C244-prakasincu}



\begin{quote}
Before the new one: say the Telugu for to transfer, then the Telugu for to blow.
\end{quote}

\subsection*{You'll want to know: \te{ప్రకాశించు}}
\textbf{\te{ప్రకాశించు}} — \emph{prakāśiñcu} — \textquotedblleft{}to shine\textquotedblright{}.

\begin{grammarlens}[title={What you've built}]
One more word for this chapter.
\end{grammarlens}

\subsection*{Your turn}
\begin{itemize}
\raggedright
  \item \emph{Say it:} \emph{prakāśiñcu}
  \item \emph{Say it:} \emph{prakāśiñcu}, once more
  \item \emph{Say it:} say \emph{prakāśiñcu}
  \item \emph{Recall:} say the Telugu for to transfer, then the Telugu for to blow, then say \emph{prakāśiñcu} again
\end{itemize}

\begin{tcolorbox}[breakable,colback=teal!4,colframe=teal!35!black,title={Before you move on}]
What does \te{ప్రకాశించు} mean? (\textquotedblleft{}To shine\textquotedblright{}.) Say it once more.
\end{tcolorbox}
\section[\texorpdfstring{\te{గడ్డకట్టు} (gaḍḍakaṭṭu)}{గడ్డకట్టు (gaḍḍakaṭṭu)}]{\te{గడ్డకట్టు} (gaḍḍakaṭṭu) — to freeze}

\label{lesson:TE-C244-gaddakattu}



\begin{quote}
Before the new one: say the Telugu for to blow, then the Telugu for to shine.
\end{quote}

\subsection*{You'll want to know: \te{గడ్డకట్టు}}
\textbf{\te{గడ్డకట్టు}} — \emph{gaḍḍakaṭṭu} — \textquotedblleft{}to freeze\textquotedblright{}.

\begin{grammarlens}[title={What you've built}]
One more word for this chapter.
\end{grammarlens}

\subsection*{Your turn}
\begin{itemize}
\raggedright
  \item \emph{Say it:} \emph{gaḍḍakaṭṭu}
  \item \emph{Say it:} \emph{gaḍḍakaṭṭu}, once more
  \item \emph{Say it:} say \emph{gaḍḍakaṭṭu}
  \item \emph{Recall:} say the Telugu for to blow, then the Telugu for to shine, then say \emph{gaḍḍakaṭṭu} again
\end{itemize}

\begin{tcolorbox}[breakable,colback=teal!4,colframe=teal!35!black,title={Before you move on}]
What does \te{గడ్డకట్టు} mean? (\textquotedblleft{}To freeze\textquotedblright{}.) Say it once more.
\end{tcolorbox}
\section[\texorpdfstring{\te{వాడిపోవు} (vāḍipōvu)}{వాడిపోవు (vāḍipōvu)}]{\te{వాడిపోవు} (vāḍipōvu) — to wilt, to fade}

\label{lesson:TE-C244-vadipovu}



\begin{quote}
Before the new one: say the Telugu for to shine, then the Telugu for to freeze.
\end{quote}

\subsection*{You'll want to know: \te{వాడిపోవు}}
\textbf{\te{వాడిపోవు}} — \emph{vāḍipōvu} — \textquotedblleft{}to wilt, to fade\textquotedblright{}.

\begin{grammarlens}[title={What you've built}]
One more word for this chapter.
\end{grammarlens}

\subsection*{Your turn}
\begin{itemize}
\raggedright
  \item \emph{Say it:} \emph{vāḍipōvu}
  \item \emph{Say it:} \emph{vāḍipōvu}, once more
  \item \emph{Say it:} say \emph{vāḍipōvu}
  \item \emph{Recall:} say the Telugu for to shine, then the Telugu for to freeze, then say \emph{vāḍipōvu} again
\end{itemize}

\begin{tcolorbox}[breakable,colback=teal!4,colframe=teal!35!black,title={Before you move on}]
What does \te{వాడిపోవు} mean? (\textquotedblleft{}To wilt, to fade\textquotedblright{}.) Say it once more.
\end{tcolorbox}
\section[\texorpdfstring{\te{మొరుగు} (morugu)}{మొరుగు (morugu)}]{\te{మొరుగు} (morugu) — to bark}

\label{lesson:TE-C244-morugu}



\begin{quote}
Before the new one: say the Telugu for to freeze, then the Telugu for to wilt.
\end{quote}

\subsection*{You'll want to know: \te{మొరుగు}}
\textbf{\te{మొరుగు}} — \emph{morugu} — \textquotedblleft{}to bark\textquotedblright{}.

\begin{grammarlens}[title={What you've built}]
One more word for this chapter.
\end{grammarlens}

\subsection*{Your turn}
\begin{itemize}
\raggedright
  \item \emph{Say it:} \emph{morugu}
  \item \emph{Say it:} \emph{morugu}, once more
  \item \emph{Say it:} say \emph{morugu}
  \item \emph{Recall:} say the Telugu for to freeze, then the Telugu for to wilt, then say \emph{morugu} again
\end{itemize}

\begin{tcolorbox}[breakable,colback=teal!4,colframe=teal!35!black,title={Before you move on}]
What does \te{మొరుగు} mean? (\textquotedblleft{}To bark\textquotedblright{}.) Say it once more.
\end{tcolorbox}
